% §2.4 Model Order
%
% Cross-reference resolution status:
%
% RESOLVED:
%   \Cref{sec:bsc}, \Cref{def:bsc}                 -- chapter 2 §2.1
%   \Cref{sec:asymmetric-channel},
%   \Cref{def:asymmetric-channel}, \Cref{def:z-channel} -- §2.2
%   \Cref{sec:channel-composition}                 -- §2.3
%   \Cref{sec:axiom-2}, \Cref{ax:block-rates-indep},
%   \Cref{cor:parameter-reduction}                 -- chapter 4
%   \Cref{ch:approximate-functions}                -- chapter 9
%
% DEFINED here:
%   \label{sec:model-order}, \label{def:model-order},
%   \label{thm:boolean-ceiling}
%
% Monograph register: definition (coarsest compatible partition; BSC is
% order 1), Boolean ceiling theorem with proof, the order-vs-parameter
% distinction with the Z-channel example and a comparison table, and the
% Axiom-2 and n-ary connections. Correctness: the partition that groups
% same-rate inputs is the COARSEST compatible one, not the "finest".

\section{Model Order}
\label{sec:model-order}

The BSC carries one crossover parameter and the asymmetric channel two, one per
input. Whether a binary channel could carry more is a question about the input
domain, answered by \emph{model order}.

\subsection*{Partitions and blocks}

A \emph{partition} of a set $S$ is a collection of non-empty, pairwise disjoint
blocks $B_1, \ldots, B_n$ whose union is $S$. In the channel model the relevant
partition is of the input domain: elements of one block share a per-block error
rate, and Axiom~2 (\Cref{sec:axiom-2}) treats the per-block rates $\{\alpha_1,
\ldots, \alpha_n\}$ as the channel's free parameters.

\begin{definition}[Model order]
\label{def:model-order}
Let $\phi : D \to D$ be a binary channel on input domain $D$. A partition
$\{B_1, \ldots, B_n\}$ of $D$ is \emph{compatible} with $\phi$ when every
element of a block shares one error probability,
\[
  \Pr\{\phi(x) \neq x\} = \alpha_i \quad \text{for all } x \in B_i.
\]
The \emph{model order} of $\phi$ is the number of blocks in the \emph{coarsest}
compatible partition: distinct error probabilities are placed in distinct
blocks, and elements are grouped only when they genuinely share a rate. Model
order is a property of the channel, the number of distinct error regimes its
input domain carries, independent of any particular parameterization.
\end{definition}

\subsection*{The Boolean ceiling}

The Boolean domain $D = \{\top, \bot\}$ admits exactly two partitions: the
\emph{trivial} one, $\{\{\top, \bot\}\}$, with both inputs in a single block
sharing one rate; and the \emph{discrete} one, $\{\{\top\}, \{\bot\}\}$, with
each input carrying its own rate.

\begin{theorem}[Boolean ceiling]
\label{thm:boolean-ceiling}
Any binary channel with input domain $\{\top, \bot\}$ has model order at most
$2$.
\end{theorem}

\begin{proof}
A partition of a set of cardinality $k$ has at most $k$ blocks, since each block
is non-empty and the blocks are disjoint. For $D = \{\top, \bot\}$ the only
partitions are the trivial (one block) and the discrete (two blocks), so the
coarsest compatible partition has one or two blocks.
\end{proof}

The BSC realizes order 1: both inputs share the crossover $\fprate$, so the
trivial partition is compatible and is the coarsest such, with channel matrix
$Q_{\mathrm{BSC}}$ from \Cref{def:bsc}. The asymmetric channel of
\Cref{def:asymmetric-channel} realizes order 2: with $\fprate \neq \fnrate$ the
two inputs cannot share a block. The ceiling is tight (order 2 is attained) but
not forced (the BSC sits at order 1). Serial composition stays within it:
composing binary channels (\Cref{sec:channel-composition}) yields a binary
channel, still of model order at most 2.

\subsection*{Model order versus parameter count}

Model order is not the same as the number of free parameters, and conflating
them is a common error. Model order counts partition blocks; parameter count
counts the real numbers needed to specify the channel. For a single Boolean
query an $n$-th order channel has $n$ free parameters, one per block, so the two
agree, but they part company in two ways. First, a constraint can reduce the
parameter count below the order: the Z-channel of \Cref{def:z-channel} has model
order 2 (its blocks carry the distinct rates $0$ and $\fprate$) yet only one
free parameter, since $\fnrate = 0$ is fixed. Second, in the multi-query setting
the parameter count grows with the number of queries while the model order does
not: a general order-$n$ channel on $m$ joint inputs carries $n(2^m - 1)$
parameters in the per-source-block convention, which the independence axioms
collapse to $nm$ by Kronecker factorization (\Cref{cor:parameter-reduction}),
all while the model order stays $n$.

\begin{center}
\begin{tabular}{l|c|c|c}
  Channel & $\alpha_1$ ($B_1 = \{\top\}$) & $\alpha_2$ ($B_2 = \{\bot\}$) & order \\
  \hline
  BSC             & $\fprate$ & $\fprate$ ($= \alpha_1$) & 1 \\
  Z-channel       & $0$       & $\fprate \neq 0$          & 2 \\
  Asymmetric      & $\fnrate$ & $\fprate$ (free)          & 2 \\
  Perfect channel & $0$       & $0$                        & 1 \\
\end{tabular}
\end{center}

The BSC and the perfect channel merge their two blocks (equal rates) and sit at
order 1; the Z-channel and the asymmetric channel keep distinct rates and sit at
order 2. The Z-channel row is the clean separation of the two counts: order 2,
one parameter.

\subsection*{Axiom 2 and the n-ary ceiling}

The partition-block structure is exactly what Axiom~2
(\Cref{ax:block-rates-indep}) acts on: it asserts that the per-block rates are
conditionally independent given the objective set, and the model order is the
dimension of that independence claim. In the order-2 Boolean case this is the
independence of $\fprate$ and $\fnrate$. The Boolean ceiling bounds the claim at
two factors, which is what keeps the framework tractable: every Boolean channel
is at most two parameters per query, and the Kronecker factorization reduces the
$m$-query count to $2m$. The machinery is not special to Booleans. For an input
domain of cardinality $k$ the same partition argument gives a ceiling of $k$,
the subject of \Cref{ch:approximate-functions}; Part~I develops the order-2
Boolean case, from which the general theory is the natural lift.
